% Propietario conservado: /Users/ruben/Documents/ChatGPT/jueces y controles/REVISION_ARGUMENTAL_APERTURAS_CIERRES_20260915/II/manuscript_es/sections/definicion_boltzmann_rev06.tex
\subsection{Definición estructural y función de la constante de Boltzmann}
\label{ii:sec:ii-definicion-boltzmann}

\begin{definition}[Transductor térmico de la energía por decisión]
Fijada una sección térmica positiva
\(\Theta_{{\rm clk},a}\in\mathcal L_\Theta\), el transductor
de Boltzmann es la aplicación lineal positiva
\(k_B:\mathcal L_\Theta\to\mathcal L_E\) que realiza
\begin{equation}
 k_B(\Theta_{{\rm clk},a})
 =\frac{\epsilon_a}{\log3},
 \qquad \epsilon_a=\frac{h_a}{108t_0}.
 \label{ii:eq:ii-definicion-boltzmann-transductor}
\end{equation}
La acción y la duración determinan \(\epsilon_a\); la decisión
trítica uniforme determina su normalización informacional.
\end{definition}

La unicidad, demostrada en la sección~\ref{ii:sec:ii-kb-aut-desarrollo},
resulta de prescribir la imagen de una sección no nula de una recta.
Las ecuaciones \eqref{ii:eq:ii-kb-aut-coordenada} y
\eqref{ii:eq:ii-kb-aut-orientaciones} precisan, respectivamente, su
coordenada en bases declaradas y la compatibilidad de un mismo
transductor con ambas orientaciones. La determinación del producto
\(k_B\Theta_{{\rm clk},a}\) precede a la individualización:
esta última conserva como datos explícitos la sección térmica, las
bases y, al reivindicar selección interna, su lector constructor.
La presión de \eqref{ii:eq:ii-kb-aut-presion}, bajo las hipótesis
allí enumeradas, normaliza el parámetro energético
\(\beta=(k_B\Theta)^{-1}\).

La función física de \(k_B\) consiste en publicar dimensionalmente
la información de una distribución: \(S=k_Bs\). Evaluar esa
entropía en una temperatura produce la energía \(\Theta S\).
Una terna uniforme aporta \(\log3\) nats; una distribución de
continuaciones aporta su entropía ponderada; una reducción cuántica
aporta la entropía de la bipartición especificada. El transductor
actúa sobre esos contenidos conservando sus dominios. La memoria
genealógica determina qué información retiene cada lectura; el
balance termodinámico declara qué operación la intercambia o la
borra. Esta separación explica la relación entre reversibilidad,
entrelazamiento y energía de decisión.
